% B卷 Q2 — 简答题
\begin{question}{0.1}
(20分) 简答题
\begin{enumerate}
\item 简述位移法的基本思路及其与力法的主要区别。
\item 简述图乘法的适用条件及计算步骤。
\end{enumerate}
\end{question}

\begin{solution}
\textbf{1. 位移法的基本思路及其与力法的主要区别}

位移法以结点位移（转角、线位移）为基本未知量。先将结构拆分为若干单跨超静定梁（三类基本杆件），
利用转角位移方程写出杆端内力与结点位移的关系，再由结点平衡条件建立位移法方程，
解出结点位移后回代求内力、作内力图。

位移法与力法的主要区别：

\textbf{基本未知量}：位移法——结点位移（几何量）；力法——多余未知力（力的量）。

\textbf{基本结构}：位移法——单跨超静定梁的组合；力法——撤除多余约束的静定结构。

\textbf{适用范围}：位移法尤其适合超静定次数高、结点位移少的结构；力法适合超静定次数低、基本结构简单的结构。

\textbf{建立方程的出发点}：位移法——平衡方程；力法——变形协调方程。

\textbf{2. 图乘法的适用条件及计算步骤}

适用条件：1. 杆件为等截面直杆（$EI = \text{常数}$）；
2. 在所要计算的区段内，$\bar{M}$ 图和 $M_P$ 图中至少有一个为直线图形；
3. 杆轴线为直线。

计算步骤：

1. 分别作出荷载弯矩图 $M_P$ 和单位力弯矩图 $\bar{M}$；

2. 在直线图形上标出面积 $A$（若均非直线则需分段），取该面积所对应图形的形心位置；

3. 在另一图形上取相应形心处的纵距 $y_c$；

4. 按公式 $\delta = \sum \frac{A y_c}{EI}$ 计算位移（注意正负号：同侧为正，异侧为负）。
\end{solution}
